\documentclass[11pt,a4paper]{article}
\usepackage[margin=25mm]{geometry}
\usepackage{amsmath,amssymb,amsthm,booktabs,longtable}
\newtheorem{theorem}{Theorem}
\newtheorem{lemma}[theorem]{Lemma}
\newcommand{\Sz}{\operatorname{Sz}}
\newcommand{\Tr}{\operatorname{Tr}}
\newcommand{\F}{\mathbb F}
\title{Excluding an order-four overlap with a Suzuki group over eight elements}
\author{Human Compression project: bounded F003 continuation}
\date{}
\begin{document}
\maketitle

We give a uniform exclusion of one remaining family in the Suzuki
rotation problem. The field of the second face and the ambient field
are unrestricted apart from the necessary divisibility conditions.
The proof uses finitely many polynomial identities over $\F_8$;
it does not enumerate an ambient Suzuki group. The identities are
supplied in full, with an exact checker using only the Python standard
library. Their derivation used computer algebra, but their verification
requires only multiplication and addition of polynomials.

The classical description of Suzuki groups, their subfield subgroups,
their ovoid action, and the reduction to two proper faces are retained
from the three accepted predecessor notes. The order-four normalization
is also established in the accepted compositum note. We recall the
parts used below. The additional argument here is the combination of
symplectic linear equations, the Suzuki characteristic polynomial, and
coefficient descent through two fields with intersection $\F_2$.
This is a bounded result. Odd middle rotations and larger first-face
fields remain outside its conclusion.

\begin{theorem}\label{main}
Let $N$ be odd, and let $S=\Sz(2^N)$ in its standard four-dimensional
representation. There do not exist $a,b,c\in S$ such that
\begin{enumerate}
\item $a$ and $c$ have odd order greater than one, and $b$ has order four;
\item $ab$, $bc$ and $abc$ are involutions;
\item $F=\langle a,b\rangle$ and $E=\langle b,c\rangle$ are genuine
Suzuki subfield groups, $F\cap E=\langle b\rangle$, and at least one
of $F,E$ is isomorphic to $\Sz(8)$.
\end{enumerate}
\end{theorem}

The assertion does not require $\langle a,b,c\rangle=S$. To exchange
the faces, replace $(a,b,c)$ by $(c^{-1},b^{-1},a^{-1})$; the three
required products are the inverses of $bc$, $ab$ and $abc$ in reverse
order. It therefore suffices to consider $F\cong\Sz(8)$.

Put $k=\F_{2^N}$ and $\theta=2^{(N+1)/2}$. On $k$ we have
$(t^\theta)^\theta=t^2$. Our coordinates are
\[
 U(x,y)=\begin{pmatrix}
 1&0&0&0\\x&1&0&0\\x^{\theta+1}+y&x^\theta&1&0\\
 x^{\theta+2}+xy+y^\theta&y&x&1
 \end{pmatrix},\qquad z(t)=U(0,t),
\]
\[
 H(\lambda)=\operatorname{diag}(\lambda^{\theta+1},\lambda,
 \lambda^{-1},\lambda^{-\theta-1}),\qquad
 J=\begin{pmatrix}0&0&0&1\\0&0&1&0\\0&1&0&0\\1&0&0&0\end{pmatrix}.
\]
These are the standard coordinates used in the predecessors, from
Bray--B\"a\"arnhielm, equations (1)--(5). In particular
\[
 U(x,y)U(w,t)=U(x+w,y+t+xw^\theta),
 \quad U(x,y)^{-1}=U(x,y+x^{\theta+1}).                 \tag{1}
\]

\section*{The characteristic polynomial}

For a four by four matrix $g$, let $c_2(g)$ be the coefficient of
$X^2$ in its characteristic polynomial. In characteristic two,
\[
 c_2(g)=\sum_{1\le i<j\le4}(g_{ii}g_{jj}+g_{ij}g_{ji}). \tag{2}
\]

\begin{lemma}\label{characteristic}
For every $g\in S$, $c_2(g)=\Tr(g)^\theta$. In particular,
if $\Tr(g)=0$, then $g^4=1$.
\end{lemma}
\begin{proof}
The standard Bruhat decomposition is $S=UH\cup UHJ U$.
For $g=U(s,t)H(\lambda)$, the diagonal entries give
\[
 c_2(g)=\lambda^{\theta+2}+\lambda^\theta+
 \lambda^{-\theta}+\lambda^{-\theta-2}=\Tr(g)^\theta.
\]
In the other cell, conjugation by the first root-group factor reduces
the calculation to $g=H(\lambda)JU(s,t)$. Multiplication gives
\begin{align*}
 \Tr(g)&=\lambda^{\theta+1}(s^{\theta+2}+st+t^\theta)
                 +\lambda s^\theta,\\
 c_2(g)&=\lambda^{\theta+2}
       (s^{2\theta+2}+s^\theta t^\theta+t^2)
                 +\lambda^\theta s^2.
\end{align*}
The second expression is the $\theta$-power of the first, since
$\theta^2$ acts as squaring. These displayed coefficient calculations
are also replayed as identities of Laurent polynomials by the checker.

The standard representation is symplectic for $J$. Its characteristic
polynomial is reciprocal, with constant term one, and is therefore
\[
 X^4+\Tr(g)X^3+c_2(g)X^2+\Tr(g)X+1.
\]
If the trace vanishes, this is $(X+1)^4$. Cayley--Hamilton gives
$(g+1)^4=g^4+1=0$, as asserted.
\end{proof}

\section*{Normalization and coverage of the first face}

Suppose the configuration of Theorem~\ref{main} exists, with
$F\cong\Sz(8)$ and $E\cong\Sz(2^e)$. The accepted cyclic-overlap
lemma gives $\gcd(3,e)=1$: the centres of the point Sylow two-groups
of the two faces have an intersection of order $2^{\gcd(3,e)}$,
whereas the cyclic group $\langle b\rangle$ has only two elements
of order dividing two. Put $K=\F_8$ and $L=\F_{2^e}$ inside $k$.
Thus
\[
 K\cap L=\F_2,\qquad [KL:L]=3.                       \tag{3}
\]
The latter degree, not merely an informal separation of coefficients,
will be needed below.

Conjugate $F$ to its standard embedding and the unique fixed point of
$b$ to infinity. The diagonal torus makes its first root coordinate
one. Conjugation by $U(h,t)$ changes its second coordinate by
$h+h^\theta$. This map has kernel $\F_2$ and image the
absolute-trace-zero subspace. Hence
\[
 b=U(1,\epsilon),\qquad \epsilon\in\{0,1\},\qquad
 b^{-1}=U(1,1+\epsilon),\qquad b^2=z(1).              \tag{4}
\]
This normalization is within $F$. On $K$ the automorphism $\theta$
is the fourth-power map.

Let $u=ab$ and $v=bc$. Neither involution fixes infinity, since
otherwise its face group fixes infinity. The retained root-group
description of the ovoid gives
\[
 u=U(x,h)^{-1}Jz(r)JU(x,h),\qquad x,h\in K,\quad r\in K^\times.
\]
Conjugation by a central root element removes $h$ and preserves $b$.
We consequently take
\[
 u=U(x,0)^{-1}Jz(r)JU(x,0).                          \tag{5}
\]
The same coordinate calculation, with arbitrary second coordinate,
gives
\[
 \Tr(b^2u)=r^\theta,\qquad
 \Tr(bu)=r^\theta(x^2+x+1)+r.                       \tag{6}
\]
Conjugation by $b$ changes the first fixed-point coordinate by one;
a central root element restores its second coordinate to zero.
The parameter $r$ remains fixed, as follows also from the first
trace in (6). We may therefore choose one of $x,x+1$.

Write $K=\F_2[z]/(z^3+z+1)$ and encode
$a_0+a_1z+a_2z^2$ by the integer $a_0+2a_1+4a_2$.
Choose $x\in\{0,2,4,6\}$ and $r\in\{1,\ldots,7\}$.
Since $bu=a^{-1}$ has odd order greater than one,
Lemma~\ref{characteristic} excludes the pairs with
$r^4(x^2+x+1)+r=0$. Exactly twenty-four pairs remain.
The field automorphism of squaring preserves $b$ and permutes them,
after choosing the representative of $x+\F_2$. Its action is
\[
 (x,r)\longmapsto (\min\{x^2,x^2+1\},r^2),           \tag{7}
\]
where the minimum refers only to the integer encoding.
The eight orbits are listed in Table~\ref{orbits}. Every admissible
first face is thus represented by a row, for each of the two values
of $\epsilon$. Including a pair which generates a smaller group
would only strengthen the exclusion; no generation test is used
to discard pairs.

\begin{table}[ht]
\centering
\begin{tabular}{lll}
\toprule
Representative & Its orbit under squaring & Certificate target\\
\midrule
$(0,2)$ & $(0,2),(0,4),(0,6)$ & $1$\\
$(0,3)$ & $(0,3),(0,5),(0,7)$ & $1$\\
$(2,1)$ & $(2,1),(4,1),(6,1)$ & $1$\\
$(2,2)$ & $(2,2),(4,4),(6,6)$ & $Q_1^3+Q_1^2+1$\\
$(2,4)$ & $(2,4),(4,6),(6,2)$ & $1$\\
$(2,5)$ & $(2,5),(4,7),(6,3)$ & $\sum_{i=0}^6Q_1^i$\\
$(2,6)$ & $(2,6),(4,2),(6,4)$ & $1$\\
$(2,7)$ & $(2,7),(4,3),(6,5)$ & $1$\\
\bottomrule
\end{tabular}
\caption{Complete orbit list. The target is the same for
$\epsilon=0$ and $\epsilon=1$; the certificates treat them separately.}
\label{orbits}
\end{table}

\section*{The linear equations and coefficient descent}

The product $abc=ub^{-1}v$ is an involution. As $u$ and $v$ are
involutions, this is equivalent to
\[
 ub^{-1}v=vbu.                                        \tag{8}
\]
Set
\[
 B=\Tr(b^2v),\qquad Q=\Tr(bv),\qquad
 B_1=B^\theta,\qquad Q_1=Q^\theta.                    \tag{9}
\]
All four elements belong to $L$. Indeed $b,v\in E$, and simultaneous
conjugation puts their matrices over $L$; traces do not change under
conjugation. The field $L$ is preserved by $\theta$.
Moreover $Q=\Tr(c)$: the traces of a symplectic matrix and its inverse
agree, so $\Tr(bv)=\Tr(vb^{-1})=\Tr(b^{-1}v)$.
The odd order of $c$ therefore gives $Q\ne0$.
To see $B\ne0$, write the involution $v$ at its fixed point away
from infinity as $U(y,t)^{-1}Jz(s)JU(y,t)$, where $s\ne0$.
The first formula of (6), which is independent of the second
coordinate, gives $B=s^\theta\ne0$.

Every symplectic involution is self-adjoint for $J$:
$v^{\mathsf T}J=Jv$. It has the shape
\[
 V=\begin{pmatrix}
 v_0&v_1&v_2&v_3\\v_4&v_5&v_6&v_2\\
 v_7&v_8&v_5&v_1\\v_9&v_7&v_4&v_0
 \end{pmatrix}.                                      \tag{10}
\]
For each of the sixteen choices of $(\epsilon,x,r)$, solve the
sixteen linear equations (8) together with the two trace equations
in (9), using the ten entries in (10) as unknowns. The coefficient
matrix over $K$ has rank eight. Its solutions are therefore
\[
             V=V(B,Q,s_0,t_0),                       \tag{11}
\]
with entries linear over $K$ and two free coordinates $s_0,t_0$.
No restriction on their field is imposed. The file
\texttt{model.json} for each case lists every entry of (11) and the
chosen free positions. The checker verifies the rank, every linear
equation, and that these positions are exactly $s_0,t_0$.
Thus (11) contains every possible $v$, rather than a selected
component of the linear solution set.

Let $I_1,\ldots,I_m\in K[s_0,t_0,B,Q,B_1,Q_1]$ be, in order,
the nonzero entries of $V^2-1$, followed by
\[
 c_2(b^2V)-B_1,\qquad c_2(bV)-Q_1.                    \tag{12}
\]
All vanish at the hypothetical configuration by
Lemma~\ref{characteristic}. The coefficient files include these
generators and polynomials $E_j\in K[B,Q,B_1,Q_1]$, together with
the complete identities
\[
                 E_j=\sum_{i=1}^m C_{ij}I_i.         \tag{13}
\]
In particular each $E_j$ vanishes there. Write
\[
                 E_j=E_{j0}+zE_{j1}+z^2E_{j2},
        \quad E_{jh}\in\F_2[B,Q,B_1,Q_1].            \tag{14}
\]
The evaluation of each $E_{jh}$ lies in $L$. Equation (3) implies
that $1,z,z^2$ are linearly independent over $L$. Hence every
$E_{jh}$ vanishes separately. This is precisely the point at which
the two face fields enter the argument.

List the distinct nonzero polynomials in (14) as
$G_1,\ldots,G_{n-1}$ and append
\[
                 G_n=wBQ-1.                         \tag{15}
\]
All these polynomials vanish on the assumed configuration by setting
$w=(BQ)^{-1}\in L$. For twelve cases the supplied identities are
\[
                 1=\sum_{i=1}^n H_iG_i,             \tag{16}
\]
which is impossible. For the other four cases they are
\[
                 p(Q_1)=\sum_{i=1}^n H_iG_i,         \tag{17}
\]
with $p=p_3=X^3+X^2+1$ or $p=p_6=\sum_{i=0}^6X^i$ as indicated
in the table. In these cases use the exact identities
\begin{align*}
 (X^5+X^4+X^3+X)p_3&=X^8+X,&
 p_3+X(X^2+X)&=1,\\
 (X^2+X)p_6&=X^8+X,&
 p_6+(X^4+X^2+1)(X^2+X)&=1.
\end{align*}
Thus a root of either $p$ in $L$ also satisfies $X^8+X=0$ and
belongs to $K\cap L=\F_2$. It then satisfies $X^2+X=0$, contradicting
the corresponding identity on the right. This excludes all sixteen
cases and proves Theorem~\ref{main}.

\section*{Certificate specification and scope}

The proof uses only the identities (13), (16) and (17), the recorded
linear substitutions, and the orbit list. It does not rely on a
software assertion that an ideal has a particular dimension or that
a computed basis is a Gr\"obner basis. For each case the accompanying
directory contains the following mathematical data.
\begin{enumerate}
\item \texttt{model.json}: the matrices $b,u$, the entries of (11),
the free positions, and the orbit represented.
\item \texttt{elimination\_certificate.txt}: all $I_i$, all $E_j$,
and all multipliers $C_{ij}$ in (13), including zero multipliers.
\item \texttt{coefficient\_mapping.json}: the correspondence between
the three coefficients of each $E_j$ and the distinct $G_i$.
\item \texttt{unit\_certificate.txt}: all $G_i$, including (15),
and all $H_i$ for (16) when the target is one.
\item \texttt{residual\_certificate.txt}, in the four other cases:
the target $p(Q_1)$ and all $H_i$ in (17).
\end{enumerate}
In these files the variable names are \texttt{ss}, \texttt{tt},
\texttt{B}, \texttt{Q}, \texttt{B1}, \texttt{Q1}, and \texttt{zz}
for $s_0,t_0,B,Q,B_1,Q_1,w$, respectively. Coefficients in $K$ are
written as polynomials in \texttt{z} reduced modulo $z^3+z+1$.

The checker reconstructs the matrices and generators, verifies every
identity by sparse polynomial arithmetic, recomputes all twenty-four
pairs and their eight orbits, and checks the displayed univariate
identities. Its characteristic-polynomial calculations treat the
symbols $s,s^\theta,t,t^\theta,\lambda,\lambda^\theta$ independently
as Laurent variables; they introduce no assumption about a finite
sample. The separately recorded negative controls corrupt exact
coefficients or omit required coverage and confirm that the relevant
equalities fail. The packet contains a full digest inventory and
bindings to the written predecessors and their reviews.

The coordinate and group-structure inputs are inherited, not claimed
as new theorems. The additional theorem excludes an order-four overlap
with one face isomorphic to $\Sz(8)$. It does not exclude odd middle rotations,
does not cover the other coprime first-face degrees, and does not
recover the whole rank theorem of the original F003 manuscript.
No new original, conjecture, Lean, or compression credit is assigned.
Distinct review of this author packet is still required.

\begin{thebibliography}{9}
\bibitem{BB}
J. N. Bray and H. B\"a\"arnhielm,
\emph{A new method for recognising Suzuki groups}, arXiv:1706.05697v1
(2017), equations (1)--(5). The exact source is pinned in the dependency
ledger.
\bibitem{Spiga}
P. Spiga, \emph{On the rank of Suzuki polytopes: an answer to Hubard and Leemans},
arXiv:2010.04080v1 (2020). The title and source edition are recorded
in the retained F003 source ledger; the present theorem is only a
bounded continuation of its recovery.
\end{thebibliography}
\end{document}
